\documentclass[11pt]{article}

\usepackage[T1]{fontenc}
\usepackage[utf8]{inputenc}
\usepackage{lmodern}
\usepackage{amsmath,amssymb,amsthm,mathtools}
\usepackage{booktabs}
\usepackage{array}
\usepackage{geometry}
\usepackage{enumitem}
\usepackage[hidelinks]{hyperref}
\hypersetup{
  pdftitle={Square Elliptic Completions and APS Boundary Information on Round Spherical Quotients},
  pdfauthor={Luqman Mahmood},
  pdfsubject={Full elliptic completion, signed boundary spectrum and phase classification}
}

\geometry{margin=1in}
\setlist{nosep}
\setlength{\parindent}{0pt}
\setlength{\parskip}{0.55em}

\newtheorem{theorem}{Theorem}[section]
\newtheorem{proposition}[theorem]{Proposition}
\newtheorem{corollary}[theorem]{Corollary}
\theoremstyle{definition}
\newtheorem{definition}[theorem]{Definition}
\newtheorem{remark}[theorem]{Remark}

\newcommand{\etaInv}{\eta}
\newcommand{\Hom}{\operatorname{Hom}}
\newcommand{\sgn}{\operatorname{sgn}}
\newcommand{\Spin}{\operatorname{Spin}}
\newcommand{\Tr}{\operatorname{Tr}}
\newcommand{\TT}{\mathrm{TT}}
\newcommand{\aux}{\mathrm{aux}}
\newcommand{\full}{\mathrm{full}}
\newcommand{\Z}{\mathbb{Z}}
\newcommand{\Q}{\mathbb{Q}}

\title{Square Elliptic Completions and APS Boundary Information\\
on Round Spherical Quotients}
\author{Luqman Mahmood\thanks{Copyright 2026 Luqman Mahmood. This paper is licensed under Creative Commons Attribution 4.0 International (CC BY 4.0); see \url{https://creativecommons.org/licenses/by/4.0/}. Software is excluded from this grant.}}
\date{30 September 2026}

\begin{document}
\maketitle

\begin{abstract}
The Atiyah--Patodi--Singer correction combines spectral asymmetry with zero
modes. We distinguish a sparse elliptic completion (one permitted off-diagonal
block is set to zero) from its fully coupled
twisted-Dirac completion, and derive the latter's signed round-boundary
spectrum from its homogeneous connection. For its round-lens character
transform we derive a closed Dedekind-sum phase law at
every rank and modulus, prove that rank one generates the entire
exponentiated ladder, and give a necessary-and-sufficient criterion for
equality under a coherent character relabeling.  We also prove infinite
families in which two adjacent ranks lose all dependence on the lens
parameter, including composite moduli.  The 23,688 exact table
entries remain reproducible by two routes. The sparse and Dirac completions
are elliptically homotopic, but no equality of their boundary spectra is
asserted. We supply neither a numerical index on a chosen filling nor a
physical interpretation of the phase.
\end{abstract}

\paragraph{Scope.}
This paper supplies an operator-level derivation, not an identification
inferred from a character match. Sections~2--4 specify the full completion,
its connection, chirality, collar and signed spectrum; the subsequent phase
proofs use the resulting transform. No
chosen-filling index, off-round extension or physical phase is claimed.

\section{Introduction}

Elliptic boundary problems naturally separate an interior symbol question
from a boundary-information question.  The first asks whether an
overdetermined differential operator can be completed to a square elliptic
one.  The second asks how much geometry survives in its spectral boundary
correction.  For the operators considered here these questions meet in a
particularly explicit form: adding the missing cross-block to a sparse
Stein--Weiss completion and normalizing the columns gives a twisted-Dirac
operator, whose round boundary operator contains the
tensor-curl family as one reducing summand.  The remaining summand is not
spectrally negligible; its eta and kernel contributions are essential to the
APS correction.

Let $(a,b)_{\rm str}$ denote the irreducible complex structure-group
$\Spin(4)$ representation
$V_a\boxtimes V_b$, where $V_m$ is the $(m+1)$-dimensional irreducible
$SU(2)$ representation.  The integer $J\geq 1$ is the tensor rank in the
normalization used below. Unsubscripted pairs in Section~2 are structure
labels; pairs describing boundary eigenspaces from Section~3 onwards are
isometry labels. Their relation is stated in Section~\ref{sec:dictionary}.
Our objects are a sparse square elliptic operator $\mathcal S_J$, a full
completion $\mathcal G_J$, its Clifford normalization $D^+_{T_J}$, the
latter's intrinsic self-adjoint boundary operator $B_J$,
and the tensor-curl block $A_J$ on a free round quotient
$Y=S^3/\Gamma$.  Flat twists are denoted by $\rho$.

The argument has three stages. Sections~2--3 identify the full elliptic
completion and derive the first-order round operator and reducing split.
Sections~4--5 assemble the complementary eta
and kernel, derive a cumulative formula with the companion tensor-curl and
McKay identities, and evaluate a finite lens-space reference table by two
independent exact routes.  Sections~6--7 use the resulting class transforms to
compare real eta data, reduced APS corrections, and their exponentials.

Generalized-gradient ellipticity and higher-spin Dirac factorization on
constant-curvature spaces are established machinery
\cite[arXiv v1, p.~1 and Thms.~3.12--3.13]{Pilca2011}
\cite[Prop.~1 and Thms.~5, 7, 10]{HommaTomihisa2021}.  Explicit spin-Dirac
spectra and twisted eta formulas on round spherical space forms are likewise
classical
\cite[Thms.~1--2]{Bar1996}
\cite[Thm.~4.3]{BoldtLauret2017}
\cite[Lemmas~2.1--2.2 and Thm.~2.5]{Gilkey1984}
\cite{CisnerosMolina2001,CacicMarcolliTeh2014}.  We use those frameworks at
the point where the exact normalization-specific calculation enters.

The principal statements are as follows.

\begin{enumerate}[label=(\roman*)]
\item The specified overdetermined Stein--Weiss operator has an explicit
square elliptic completion for every integer $J\geq1$, with a closed positive
determinant-square law.
\item On free round spherical quotients,
\[
 B_J\cong A_J\oplus B_J^{\aux}
\]
is an orthogonal reducing split.  The eta invariant and kernel dimension of
$B_J^{\aux}$ have exact all-rank representation-ring formulas.
\item An exact calculation covers every $2\leq p\leq24$, every unit
$q$, both spin lifts where distinct, $J=1,\ldots,6$, and every flat character.
It contains $23{,}688$ exact triples $(\etaInv,h,\xi)$, with independent
rank-seven and rank-eight checks.
\item The complete class transforms of $\etaInv$, $h$, and
$\xi=(\etaInv+h)/2$ are rational character expressions.  Reduced ranks one
and two have universal nonidentity-class support, while the rank-one
reconstruction theorem itself is known.
\item Exact finite computation refutes universal oriented adjacent-rank
reconstruction.  Sixteen rows represent four colliding oriented-orbit
pairs; two collisions remain nonhomeomorphic even after orientation reversal.
\item For the derived transform, a Dedekind-sum formula gives every
phase, rank one generates every higher phase, and a complete coherent-label
collapse criterion follows. The pair $L(16,1)$ and $L(16,9)$ is a corollary.
Infinite blind adjacent-rank families include admissible composite moduli.
\end{enumerate}

The arithmetic uses classical Fourier--Dedekind sums and congruences; no
priority claim follows from the local verification. The transform-specific
deductions are separated from those inputs and from any filling-index claim.

\section{The square completion}

\subsection{A sparse elliptic square}

The generalized-gradient construction begins with the Levi-Civita covariant
derivative and orthogonal projection onto irreducible $\Spin(4)$ summands.
Pilca's ellipticity classification provides the ambient framework
\cite[arXiv v1, p.~9, Thms.~3.12--3.13]{Pilca2011}.  The essential data here are
the exact representation choice, column normalization, and determinant law.

For $J\geq2$, choose the right-handed mirror of the sparse completion and set
\[
 \mathcal S_J:
 (J,J)\oplus(J,J-2)
 \longrightarrow
 (J+1,J-1)\oplus(J-1,J-1).
\]
Write $E=(J,J)$, $C=(J,J-2)$, $F_1=(J+1,J-1)$ and $F_2=(J-1,J-1)$.
If $P_{UV}\nabla$ denotes covariant differentiation followed by the
orthogonal projection from $U$ to $V$, the sparse operator is
\[
 \mathcal S_J=
 \begin{pmatrix}P_{EF_1}\nabla&0\\P_{EF_2}\nabla&P_{CF_2}\nabla\end{pmatrix}.
\]
There are three nonzero entries, not four. In a
left-handed presentation the auxiliary summand is $(J-2,J)$; the two choices
are mirrors, not distinct ellipticity statements.  The source and target both
have complex rank $2J(J+1)$.  At $J=1$ the negative-label auxiliary summand is
absent and the same statement is read in the reduced two-summand convention.

Fix a unit covector.  Restriction to its diagonal $SU(2)$ stabilizer makes the
symbol block triangular.  With orthonormal irreducible bases, the diagonal
amplitude squares are
\begin{align}
 \alpha_r^2&=\frac{r(r+1)}{2(J+1)^2}, &&1\leq r\leq J,\\
 \beta_r^2&=\frac{(J-r)(J+r+1)}{2(J+1)^2}, &&0\leq r\leq J-1,\\
 \gamma_r^2&=\frac{r(r+1)}{2(J^2-1)}, &&1\leq r\leq J-1.
\end{align}
Every factor in its actual range is positive.  Stabilizer transitivity and
symbol homogeneity then extend invertibility from the chosen unit covector to
every nonzero covector.

\begin{theorem}[Exact square completion]
For every integer $J\geq1$, the operator $\mathcal S_J$ is square elliptic.
For $J\geq2$, its invariant determinant-square is
\[
 \left(\frac{J}{2(J+1)}\right)^{2J+2}
 \prod_{r=1}^{J-1}
 \left(
  \frac{r^2(r+1)^2}{4(J+1)^2(J^2-1)}
 \right)^{2r+1}>0.
\]
\end{theorem}

\begin{proof}
At a unit covector, restriction to the diagonal $SU(2)$ stabilizer makes the
symbol block triangular. Its one-copy endpoints are $\beta_0$ and
$\alpha_J$; each intermediate two-copy determinant is $\alpha_r\gamma_r$.
All are nonzero. These amplitude formulas can be obtained from the
spin-half projectors: for spin $j$, the projections in
$\tfrac12\otimes j$ are $(2\boldsymbol S\cdot\boldsymbol J+j+1)/(2j+1)$
and $(j-2\boldsymbol S\cdot\boldsymbol J)/(2j+1)$. At a unit covector
the vector representation is the normalized singlet of its two spin halves.
Thus $\langle2\boldsymbol S_L\cdot\boldsymbol J_L\rangle=
\langle2\boldsymbol S_R\cdot\boldsymbol J_R\rangle=0$ and their product
has expectation $-\boldsymbol J_L\cdot\boldsymbol J_R$.
Substitute $2\boldsymbol J_L\cdot\boldsymbol J_R=
r(r+1)-j_L(j_L+1)-j_R(j_R+1)$ into the two projectors to obtain the displayed
squares. The symbol is therefore invertible at that
covector.  Stabilizer transitivity and homogeneity give invertibility at every
nonzero covector, proving square ellipticity.  Multiplying the squared
diagonal factors with their stabilizer multiplicities $2r+1$ gives the stated
determinant-square formula.  For $J=1$ the absent auxiliary column is omitted
and the same argument applies to the reduced block.
\end{proof}

This formula belongs to the stated sparse normalization.  It is not detached
from the operator and advertised as a general determinant construction.

\subsection{The full completion and its twisted-Dirac identification}

Define the half-spinors by
\[
 S^+=(0,1)_{\rm str},\qquad S^-=(1,0)_{\rm str},\qquad
 \operatorname{Sym}^2 S^\pm=\Lambda^2_\pm.
\]
These definitions incorporate the structure/isometry dictionary of
Section~\ref{sec:dictionary}; the superscripts are not the order of the
structure factors. Let
$T_J=\operatorname{Sym}^J S^-\otimes\operatorname{Sym}^{J-1}S^+=(J,J-1)_{\rm str}$.
The chiral twisted-Dirac operator
\[
 D^+_{T_J}:S^+\otimes T_J\longrightarrow S^-\otimes T_J
\]
has source and target decomposition
\[
 (J,J)\oplus(J,J-2)
 \longrightarrow
 (J+1,J-1)\oplus(J-1,J-1),
\]
and uses all four projected gradients. Choose the constant unitary
identifications of the irreducible summands and their phases so that
\[
 \mathcal G_J=
 \begin{pmatrix}P_{EF_1}\nabla&P_{CF_1}\nabla\\
 P_{EF_2}\nabla&P_{CF_2}\nabla\end{pmatrix}.
\]
The omitted entry of $\mathcal S_J$ is essential. Multiplying the two
source columns of $\mathcal G_J$ by
\[
 c_E=\sqrt{\frac{2(J+1)}{J}},
 \qquad
 c_C=\sqrt{\frac{2(J-1)}{J}},
\]
with the second column absent at $J=1$, gives exactly the Clifford symbol.
In intermediate diagonal spin $1\leq r<J$, the
recoupling matrix is
\[
 R_{J,r}=
 \begin{pmatrix}
 \sqrt{t_r}&-\sqrt{1-t_r}\\
 \sqrt{1-t_r}&\sqrt{t_r}
 \end{pmatrix},
 \qquad
 t_r=\frac{r(r+1)}{J(J+1)}.
\]
The column norms follow from the projector formulas above; orthogonality
then fixes the second column as $(-\sqrt{1-t_r},\sqrt{t_r})^T$.
At spins $0,J$ the normalized one-copy symbol is an isometry.
Equivariance gives the Clifford symbol at every covector. Moreover, this is
an equality of operators, not just symbols:
\[
 D^+_{T_J}=\mathcal G_J\operatorname{diag}(c_E,c_C).
\]
Both sides are the same parallel bundle map composed with the same
Levi-Civita-induced connection. There is no adjustable zeroth-order term.
The four-block framework is classical
\cite[Sec.~2, Eq.~(2.3)]{HommaTomihisa2021}; the normalization is fixed here.

For clarity, the determinant-square of the unnormalized full completion is
\[
 |\det\sigma(\mathcal G_J)|^2=
 \left(\frac{J}{2(J+1)}\right)^{(J+1)^2}
 \left(\frac{J}{2(J-1)}\right)^{J^2-1},\qquad J\geq2,
\]
whereas the normalized Clifford determinant-square is $1$ at a unit
covector. These are not the sparse determinant. At $J=2$ the sparse and
full values are respectively $1/14348907$ and $1/19683$.

There is nevertheless an elliptic homotopy: in a normalized intermediate
block multiply the upper-right entry of $R_{J,r}$ by $s\in[0,1]$.
Its determinant is $t_r+s(1-t_r)>0$, and the endpoints are unchanged.
This preserves the interior symbol class. It does not identify the sparse
operator's boundary spectrum with that of $D^+_{T_J}$. From now on $B_J$
always refers to the fully coupled Dirac completion.

Although $S^\pm$ and $T_J$ individually require spin lifts, their tensor
products above descend to $SO(4)$: the diagonal central element
$(-1,-1)\in\Spin(4)$ acts by
$(-1)(-1)=1$. The operator therefore exists on an oriented Riemannian
four-manifold without a chosen global spin structure, using these associated
bundles. Spin notation gives its local factorization. On the boundary a
change of spin lift acts on both factors and cancels, explaining the
spin-lift independence of the tensorial spectrum.

\section{The round boundary operator}

\subsection{Structure labels versus isometry labels}
\label{sec:dictionary}
On $\mathbb H$ oriented by $(1,i,j,k)$ use the isometry action
$q\mapsto zqw^{-1}$, with the left factor first. Left infinitesimal
multiplications correspond to the self-dual two-forms
$e^{01}+e^{23}$ and their cyclic partners; right multiplications correspond
to the anti-self-dual forms $e^{01}-e^{23}$ and their partners.
For the $d^-$ choice below, Section~2's target $(2,0)_{\rm str}$ is
$\Lambda^2_-$, so the explicit dictionary is
\[
 (a,b)_{\rm str}=(b,a)_{\rm iso}.
\]
These distinguish structure-group labels from isometry labels. For fundamental
isometry characters $x=\chi_1(z)$, $y=\chi_1(w)$ and $P=x-y$, the
source-minus-target character of the printed symbol is
\[
 -P_{\rm str}\chi_J^{\rm str,L}\chi_{J-1}^{\rm str,R}
       =P\chi_{J-1}(z)\chi_J(w).
\]
The identity follows directly from
$\chi_{r+1}+\chi_{r-1}=\chi_1\chi_r$; it is valid at every rank.
At $J=1$, the independent exterior-algebra calculation reads
\[
 \frac{\operatorname{tr}\Lambda^1-\operatorname{tr}\Lambda^2_-
       -\operatorname{tr}\Lambda^0}{\det(1-g)}
 =\frac{xy-(y^2-1)-1}{(x-y)^2}=\frac y{x-y}.
\]
Thus the swapped symbol term, rather than the same-labelled one, agrees
with the proposed fixed-point numerator. Together with
Christoffel--Darboux below, its difference from the reduced transform is
minus the kernel character. This is a character-level consistency check;
a full filling statement is intentionally not inferred from it.

\subsection{Collar, Clifford convention, and a first-order calibration}
\label{sec:collar}

Give $Y$ its boundary orientation and write the product orientation as
$dx\wedge\mathrm{vol}_Y$, with $x$ increasing outwards.  The inward collar
coordinate is $u=-x$.  We define $B_J$ by
\[
 D^+_{T_J}=c(du)(\partial_u+B_J)
\]
using the product Levi-Civita connection and the induced connection on $T_J$.
The APS condition annihilates the nonnegative spectral part of $B_J$; its
index correction is $-\xi(B_J)$, using the inward-normal convention of
\cite[Theorem~3.10, p.~57; equations~(4.1)--(4.3), pp.~60--62]{APS1975I}.
A change of the normal factorization without
changing this spectral convention is not innocuous.

Here is a calibration that fixes the sign independently of a squared
operator.  At $J=1$ the Clifford-normalized completion is
$(\sqrt{2}d^-,\delta)$ on one-forms.  Write a one-form as $f\,dx+\beta$ and
identify an anti-self-dual two-form with
$(dx\wedge v-*_Yv)/\sqrt{2}$.  In the resulting target coordinates its value is
\[
 (\partial_x\beta-d_Yf-*_Yd_Y\beta,\ -\partial_xf+\delta_Y\beta).
\]
The normal symbol sends $(f,\beta)$ to $(\beta,-f)$.  Factoring with the
inward coordinate therefore gives
\begin{equation}\label{eq:j1-collar}
 B_1=\begin{pmatrix}0&\delta_Y\\ d_Y&*_Yd_Y\end{pmatrix}.
\end{equation}
In particular, the co-closed one-form block is positive curl.  On the unit
quaternions with the boundary orientation, the right-invariant form
$dq\,q^{-1}$ has curl $+2$ and type $(2,0)$; its left-invariant counterpart
has curl $-2$ and type $(0,2)$.  Thus the positive TT branch is
$(n+2r,n)$, with eigenvalue $n+r+1$.  This is the convention of the companion
tensor-curl papers~\cite{Mahmood2026LensEta,Mahmood2026DiracToTensorCurlEtaTransfer},
not its negative.

For use at every rank, view the unit sphere as the symmetric space
$\Spin(4)/\Spin(3)$.  The canonical and Levi-Civita connections agree here;
the homogeneous Dirac formula has no cubic bracket term
\cite[Sections~2--3]{Agricola2002}. The following calculation also fixes the
twisting connection, not merely a Casimir eigenvalue.
Let $e_k\in\{i,j,k\}$ and take horizontal generators
$X_k=(e_k/2,-e_k/2)$; the curve $\exp(tX_k)\cdot1=\exp(te_k)$ has unit
speed. A section of the associated spinor-twist bundle is an equivariant
function on $\Spin(4)$, differentiated horizontally by the canonical
connection. In its Peter--Weyl decomposition write $L_k,R_k$ for the
Hermitian generators on the contragredient representation of type $(p,q)$.
The derivative along $X_k$ is $i(L_k-R_k)$. On the boundary spinor factor
put $c_Y(e_k)=2iS_k$, $S=1/2$, with the Clifford sign selected by
\eqref{eq:j1-collar}. Consequently, on the invariant multiplicity space
for a twisting summand $V_{2r-1}$, the actual first-order operator is
\begin{equation}\label{eq:signed-angular}
 D=2\boldsymbol S\mathbin{\cdot}(\boldsymbol R-\boldsymbol L),
 \qquad S=\tfrac12,\quad L=\tfrac p2,\quad R=\tfrac q2.
\end{equation}
Here the diagonal invariance condition says that $\boldsymbol S+\boldsymbol
L+\boldsymbol R$ couples to twisting spin $r-1/2$. The $SU(2)$ duals have
the same highest weights, so $p,q$ remain the isometry labels. The
horizontal bracket has no horizontal component, hence neither the spin nor
twist connection contributes an additional horizontal algebraic term.
This derives \eqref{eq:signed-angular} from the specified connection. The
first-order $J=1$ calibration fixes the remaining overall sign; it is not
inferred from the Casimir square used below.

\subsection{A reducing tensor-curl sector}

On a product collar, the tangential operator of $D^+_{T_J}$ is the intrinsic
Dirac operator on
\[
 S_Y\otimes T_J|_{\Spin(3)}
 =W_0\oplus2W_1\oplus\cdots\oplus2W_{J-1}\oplus W_J.
\]
The rank-$J$ transverse-traceless compression is the tensor-curl operator
$A_J$.  On the round three-sphere its two helicities have $\Spin(4)$ support
$(n+2J,n)$ and $(n,n+2J)$.  A lower-spin bundle $W_r$, $r<J$, only contains
types with $|p-q|\leq2r<2J$.  Equivariance therefore prohibits mixing between
the TT helicities and the auxiliary summands.

\begin{theorem}[Round reducing split]
On $S^3$ and every free round spherical quotient,
\[
 B_J\cong A_J\oplus B_J^{\aux}
\]
as an orthogonal reducing split.  Consequently, for every flat twist $\rho$,
\begin{align}
 \etaInv_\rho(B_J)&=\etaInv_\rho(A_J)+\etaInv_\rho(B_J^{\aux}),\\
 h_\rho(B_J)&=h_\rho(A_J)+h_\rho(B_J^{\aux}).
\end{align}
\end{theorem}

\begin{proof}
The two helicity branches of $A_J$ have $\Spin(4)$ types $(n+2J,n)$ and
$(n,n+2J)$.  Every lower-spin summand $W_r$, $r<J$, has
$|p-q|\leq2r<2J$, so no irreducible type occurring in the tensor-curl sector
occurs in the auxiliary sector.  Since $B_J$ is equivariant and self-adjoint,
the orthogonal projections onto these inequivalent supports commute with it.
This gives the reducing split on $S^3$; equivariance makes the decomposition
descend to a free quotient and persist after a flat twist.  Additivity of eta
and kernel over an orthogonal direct sum gives the two displayed identities.
\end{proof}

The word ``round'' is essential.  No corresponding reducing theorem is used
for Berger or arbitrary boundary metrics.

\subsection{Local index density}

Let $x_+,x_-$ be the formal half-spinor roots, with
$x_+^2=(p_1+2e)/4$ and $x_-^2=(p_1-2e)/4$.
The weight-square sum in $\operatorname{Sym}^n\mathbb C^2$ is
\[
 \sum_{k=0}^n(n-2k)^2=\frac{n(n+1)(n+2)}3.
\]
Expanding
$\mathrm{ch}(\operatorname{Sym}^J S^-)
\mathrm{ch}(\operatorname{Sym}^{J-1}S^+)$ through degree four and multiplying
by $\widehat A=1-p_1/24+\cdots$ gives, for the displayed positive-chiral
source,
\[
[\widehat A\,\mathrm{ch}(T_J)]_4
 =\frac{J(J+1)}{12}
\bigl((J^2+J-1)p_1-(2J+1)e\bigr).
\]
On a nonspin manifold the left side denotes the invariant polynomial
computed in local spin lifts; the right side is globally defined.
On a closed four-manifold this integrates to
\[
 \frac{J(J+1)}{12}
 \bigl(3(J^2+J-1)\tau-(2J+1)\chi\bigr).
\]
At $J=1$ the expression is $(\tau-\chi)/2$, matching the calibrated
$d^-+\delta$ control.  We do not choose a filling in this paper, so the local
density is not converted into a numerical APS index.

\section{Eta, kernel, and the all-rank auxiliary assembly}

\subsection{Classical quotient machinery and APS normalization}

The ordinary spin-Dirac spectrum on $S^3$ and free spherical quotients is
described by character generating functions
\cite[pp.~74--76, Thms.~1--2]{Bar1996}.  Lens-space multiplicities admit an
explicit affine-congruence-lattice formula
\cite[arXiv v1, p.~14, Thm.~4.3]{BoldtLauret2017}.  Twisted eta and defect formulas on
odd-dimensional spherical space forms, including spin-structure typing, are
classical
\cite[pp.~427--432]{Gilkey1984}.  For twisted Dirac spectra on $S^3/\Gamma$ we
cite Cisneros-Molina's research article \cite{CisnerosMolina2001} together
with the theorem-level reproduction
\cite[pp.~480--481, Thm.~2.2]{CacicMarcolliTeh2014}.  No physical inference is
taken from the latter article.

For a self-adjoint boundary operator $B$, the APS reduced correction is
\begin{equation}\label{eq:xi-definition}
 \xi_\rho(B)=\frac{\etaInv_\rho(B)+h_\rho(B)}{2},
 \qquad h_\rho(B)=\dim\ker B_\rho.
\end{equation}
The kernel term is part of the definition and of the index correction
\cite[p.~414, Eq.~(3.1) and Thm.~3.3]{APS1975}.  Eta-only quantities below are
always labeled separately.

\subsection{The auxiliary formula}

For $r\geq1$, define the finite virtual character
\[
 E_r=\sum_{m=0}^{2r-2}
 \sgn(m-r+1)
 [V_m\boxtimes V_{2r-2-m}].
\]
Let $\langle\rho,E_r\rangle$ denote the corresponding finite character
pairing.  In the angular-momentum calculation let $C_L,C_R$ be the Casimirs
of the two factors of the type $(p,q)$, with
$C(V_n)=n(n+2)/4$.  The auxiliary correction is as follows.

\begin{theorem}[All-rank auxiliary eta and kernel]
For every integer $J\geq1$, free round spherical quotient $S^3/\Gamma$, and
flat twist $\rho$,
\begin{align}
 \etaInv_\rho(B_J^{\aux})
 &=\sum_{r=1}^{J-1}\etaInv_\rho(A_r)
   -\sum_{r=1}^{J}\langle\rho,E_r\rangle,\label{eq:aux-eta}\\
 h_\rho(B_J^{\aux})
 &=\sum_{r=1}^{J}
 \dim\Hom_\Gamma
 \bigl(\rho,V_{r-1}\boxtimes V_{r-1}\bigr).\label{eq:aux-kernel}
\end{align}
Therefore
\begin{align}
 \etaInv_\rho(B_J)
 &=\sum_{r=1}^{J}\etaInv_\rho(A_r)
   -\sum_{r=1}^{J}\langle\rho,E_r\rangle,\label{eq:full-eta}\\
 h_\rho(B_J)&=h_\rho(A_J)+h_\rho(B_J^{\aux}).\label{eq:full-kernel}
\end{align}
\end{theorem}

\begin{proof}
The diagonal restriction is
$T_J|_{\Spin(3)}=\bigoplus_{r=1}^J V_{2r-1}$.  On its $r$th summand,
\eqref{eq:signed-angular} and the angular-momentum commutators give
\[
 D^2=2(C_L+C_R)+\tfrac34-C_{S+L+R}
     =2(C_L+C_R)+1-r^2.
\]
To fix the signed spectrum, couple $L+R$ first.  In a multiplicity-two
space the two intermediate spins are $t=r-1,r$.  The diagonal matrix elements
of $D$ are respectively $(C_R-C_L)/r$ and $-(C_R-C_L)/r$.
For $r=1,t=0$ the first is zero, as $C_L=C_R$ there.  Hence the trace is
zero, and the square law makes the two eigenvalues opposite.  They cancel
in eta before taking the quotient.

On the one-copy type $(n+2r,n)$ only $t=r$ is available, giving
$D=(C_L-C_R)/r=n+r+1$; the mirror type has the negative eigenvalue.
On the finite one-copy type $(m,2r-2-m)$ only $t=r-1$ is available, giving
\[
 D=(C_R-C_L)/r=r-1-m,\qquad 0\leq m\leq2r-2.
\]
Its signed character is therefore $-E_r$, not $E_r$.  The only zero type is
$V_{r-1}\boxtimes V_{r-1}$.  Passing to $\Gamma$-invariants after twisting by
$\rho$, and then summing in $r$, gives
the full formula.  Removing the $r=J$ TT ladder leaves
\eqref{eq:aux-eta}; the TT ladders have no kernel, so all zero types remain
in \eqref{eq:aux-kernel}.  The square law alone would not have fixed the
relative sign between the infinite and finite branches.
\end{proof}

\subsection{Auxiliary characters and cumulative assembly}

The generic odd-Chebyshev and representation-ring
algebra used in the cumulative calculation is developed in
\cite{Mahmood2026BivariateMcKay}.  It is algebraic input only; the square
completion, reducing split, local density, and auxiliary APS formulas are the
boundary results proved here.  The cumulative corollary below also uses the
rankwise Dirac-to-tensor-curl eta formula of
\cite{Mahmood2026DiracToTensorCurlEtaTransfer}, restricted to the common
hypotheses of the two papers.

Every multiplicity-two representation has a paired
$(+\lambda,-\lambda)$ spectrum and cancels before quotient projection.  The
unpaired infinite branches are the tensor-curl ladders; what remains is the
finite $-E_r$ contribution and the zero representation counted in
\eqref{eq:aux-kernel}.

The auxiliary eta does not vanish universally.  For the trivial sector at
ranks $J=1,2,3,4$, the exact values include
\[
 S^3/Q_8:\quad 0,\ -\frac94,\ -\frac{15}{4},\ -\frac32,
 \qquad
 S^3/(2T):\quad 0,\ -\frac{59}{36},\ -\frac{55}{12},\ -\frac{43}{6}.
\]
Here the groups act by right multiplication.  Both the signed spectral
assembly and direct class averaging reproduce these values.  Thus the
auxiliary term carries genuine nonabelian representation data.

The higher-spin operators, sphere spectra, twisted quotient formulas, eta
invariant, and APS kernel correction are established components.  What was
not found in the surveyed primary sources is the combined formula
\eqref{eq:aux-eta}--\eqref{eq:aux-kernel} in this natural round-boundary
normalization.

\subsection{Bounded cumulative eta and reduced correction}

Fix the common domain of the cited rankwise transfer and the boundary
assembly above: the same free effective round spherical quotient $S^3/\Gamma$,
injective $\Spin(4)$ lift, orientation, operator and copy-multiplicity
normalization, and flat twist.  Let $A_L,A_R$ be multiplication by the two
lifted defining representations in the common representation-ring
convention, let $U_m$ be the Chebyshev polynomial of the second kind, and put
\[
 C_J(A)=U_{J-1}(A/2)U_J(A/2),\qquad
(\boldsymbol e_r)_\rho=\langle\rho,E_r\rangle,\qquad
 \boldsymbol e_{\leq J}=\sum_{r=1}^{J}\boldsymbol e_r.
\]
Write $\boldsymbol\eta^B_J=(\etaInv_\rho(B_J))_\rho$ and let
$\boldsymbol\eta^D_C$ be the twisted spin--Dirac eta vector in the convention
of \cite{Mahmood2026DiracToTensorCurlEtaTransfer}: its positive branch is
$V_n\boxtimes V_{n+1}$.  This is the opposite of the untwisted Dirac sign
induced by the signed-angular boundary convention of the present paper.
The subscript $C$ records the companion convention; the factor $-1/2$ below
uses that sign and must not be combined with its negative.

\begin{corollary}[Cumulative boundary eta and reduced correction]
\label{cor:cumulative-boundary-eta-xi}
For every integer $J\geq1$ on the common domain above,
\begin{align}
 \boldsymbol\eta^B_J
 &=-\frac12\bigl[C_J(A_L)+C_J(A_R)\bigr]
   \boldsymbol\eta^D_C-\boldsymbol e_{\leq J},
 \label{eq:cumulative-boundary-eta}\\
 \xi_\rho(B_J)
 &=\frac12\left\{\etaInv_\rho(B_J)+h_\rho(A_J)
   +\sum_{r=1}^{J}\dim\Hom_\Gamma
   \bigl(\rho,V_{r-1}\boxtimes V_{r-1}\bigr)\right\}.
 \label{eq:cumulative-boundary-xi}
\end{align}
\end{corollary}

\begin{proof}
The rankwise transfer theorem of
\cite{Mahmood2026DiracToTensorCurlEtaTransfer} gives
$\boldsymbol\eta^{\TT}_r=-\tfrac12(A^L_{2r-1}+A^R_{2r-1})
\boldsymbol\eta^D_C$.  In the common representation-ring convention,
$A^X_m=U_m(A_X/2)$.  Substitute this formula for each
$\etaInv_\rho(A_r)$ in \eqref{eq:full-eta}, then apply the two separate
Chebyshev identities from \cite{Mahmood2026BivariateMcKay},
$\sum_{r=1}^{J}U_{2r-1}(A_X/2)=C_J(A_X)$ for $X=L,R$.  This gives
\eqref{eq:cumulative-boundary-eta}.  Equation
\eqref{eq:cumulative-boundary-xi} follows from
\eqref{eq:xi-definition}, \eqref{eq:full-kernel}, and
\eqref{eq:aux-kernel}.
\end{proof}

This corollary is a composition of two cited inputs with the boundary
assembly proved here; it is not a new operator intertwining or isospectrality
statement.  Its hypotheses are exactly the common round quotient, lift,
orientation, normalization, and flat-twist hypotheses stated above.

\section{Exact round-lens calculations}

\subsection{Domain and two routes}

We next calculate a finite-domain reference table for round lens spaces and
ranks.  The eta, kernel, and reduced correction values are independently
reconstructed by assembly and cyclotomic Fourier routes.  The calculation is
not an all-level classification or a numerical index on a chosen filling.

The calculation domain is
\[
 2\leq p\leq24,
 \quad q\in(\Z/p\Z)^\times,
 \quad J=1,\ldots,6,
\]
with both spin lifts where distinct and every flat character $a\in\Z/p\Z$.
It comprises 242 lens/lift cells, 1,452 rank rows, and 23,688 exact triples
\[
 \bigl(\etaInv_a(B_J),h_a(B_J),\xi_a(B_J)\bigr).
\]

Route A uses the all-rank split
\eqref{eq:aux-eta}--\eqref{eq:full-kernel}, the exact
Hurwitz/quasipolynomial tensor-curl formula, and direct weight enumeration for
the finite terms.  Route B starts from the closed geometric-twist class
transform in Section~\ref{sec:class-transforms}, recovers per-rank blocks by
transform differences, and applies exact finite Fourier inversion.  The
implementation works in
$\Q[X]/(1+X+\cdots+X^{p-1})$, which retains all nonidentity $p$th roots,
including imprimitive ones; the odd lift makes both character exponents
integral.  Route B uses no Hurwitz continuation.  At the identity
class the eta coefficient is zero by left-right spectral symmetry, while kernel and
finite terms are exact dimension counts.

Across the finite domain the routes agree in 94,752 exact value checks.  The
complete tables are supplied as ancillary machine-readable data, with both
route implementations and exact assertion-driven checks.  Two further
out-of-domain cells are supplied as additional reproducibility checks, not
as unobserved predictions.

\subsection{Selected values and independent checks}

Table~\ref{tab:selected-values} records a small human-readable selection.
The $L(17,3)$ row checks the auxiliary formula directly; the other rows were
calculated outside the working domain without changing the formulas.

\begin{table}[ht]
\centering
\small
\begin{tabular}{@{}lrrrrr@{}}
\toprule
space & $J$ & $a$ & $\etaInv_a(B_J)$ & $h_a(B_J)$ & $\xi_a(B_J)$\\
\midrule
$L(17,3)$  & 6 & 5  & $-9/17$  & 6 & $93/34$\\
$L(17,3)$  & 6 & 12 & $-9/17$  & 6 & $93/34$\\
$L(29,12)$ & 7 & 3  & $-13/29$ & 5 & $66/29$\\
$L(29,12)$ & 7 & 26 & $-13/29$ & 5 & $66/29$\\
$L(27,7)$  & 8 & 9  & $-1/9$   & 7 & $31/9$\\
\bottomrule
\end{tabular}
\caption{Selected exact values from the assembly and Fourier calculations.}
\label{tab:selected-values}
\end{table}

At $L(17,3)$, $J=6$, $a=5$ and its conjugate $a=12$, the auxiliary values are
$\etaInv(B_6^{\aux})=-8/17$ and $h(B_6^{\aux})=6$, while the full eta is
$-9/17$.  Conjugation, orientation reversal, spin-lift invariance, exact
denominator typing, and the identities
\[
 \sum_a\etaInv_a(B_J)=0,
 \qquad
 \sum_a h_a(B_J)=\frac{J(J+1)(2J+1)}6
\]
are verified throughout the table domain.

\section{Class transforms and oriented reconstruction}
\label{sec:class-transforms}

Let $L(p,q)$ carry the fixed odd lift $\widetilde q$.  For a nonidentity
cyclic class represented by $k$, $1\leq k\leq p-1$, put
\[
 \zeta=e^{\pi i/p},\qquad
 e_L=k(1+\widetilde q),\qquad
 e_R=k(1-\widetilde q),
\]
\[
 P_k=(\zeta^k-\zeta^{-k})
 (\zeta^{k\widetilde q}-\zeta^{-k\widetilde q}).
\]
For the $SU(2)$ character
\[
 \chi_m(z)=\sum_{j=0}^m z^{m-2j}
 =\frac{z^{m+1}-z^{-(m+1)}}{z-z^{-1}},
\]
with the polynomial value used also at $z=\pm1$, write
$\chi_m^L=\chi_m(\zeta^{e_L})$ and
$\chi_m^R=\chi_m(\zeta^{e_R})$.  On nonidentity free classes, exact Fourier
inversion of the complete flat-twist vectors gives
\begin{equation}\label{eq:eta-transform}
 \widehat\etaInv_J^{\full}(k)
 =\frac{
  \chi_J^L\chi_{J-1}^R+\chi_{J-1}^L\chi_J^R
 }{P_k}.
\end{equation}
To see the relative sign algebraically, put $x=\chi_1^L$, $y=\chi_1^R$,
so $P_k=x-y$.  The recurrence $\chi_{m+1}(x)=x\chi_m(x)-\chi_{m-1}(x)$
telescopes the finite character to
\[
 P_k\sum_{r=1}^J E_r
 =(\chi_J^L-\chi_J^R)(\chi_{J-1}^L-\chi_{J-1}^R).
\]
The positive-TT cumulative numerator is
$\chi_J^L\chi_{J-1}^L+\chi_J^R\chi_{J-1}^R$.
Subtracting the finite term gives exactly the cross numerator in
\eqref{eq:eta-transform}.

The kernel character is
\[
 H_J=\sum_{r=0}^{J-1}\chi_r^L\chi_r^R.
\]
Christoffel--Darboux gives
\[
 P_kH_J
 =\chi_J^L\chi_{J-1}^R-\chi_{J-1}^L\chi_J^R.
\]
Combining eta and kernel according to
\eqref{eq:xi-definition} selects one chirality:
\begin{equation}\label{eq:xi-transform}
 \widehat\xi_J(k)
 =\frac{\chi_J^L\chi_{J-1}^R}{P_k}.
\end{equation}

\subsection{Blind classes}

The full-eta blind set is the cyclotomic divisibility locus of
\[
 \chi_J^L\chi_{J-1}^R+\chi_{J-1}^L\chi_J^R.
\]
For every even $p$, the half-turn class $k=p/2$ is blind for every rank and
lift.  Thus no finite collection of full-eta ranks has complete Fourier
support at all levels.

The reduced transform vanishes precisely when
\[
 (J+1)e_L\equiv0\pmod p,
 \qquad e_L\not\equiv0,p\pmod{2p},
\]
or
\[
 Je_R\equiv0\pmod p,
 \qquad e_R\not\equiv0,p\pmod{2p}.
\]
Reduced ranks one and two have no common blind nonidentity class for any
$p,q$, or lift.  This is a universal support statement, not by itself a
reconstruction theorem.

\subsection{The known rank-one base}

At $J=1$, the auxiliary eta is zero and the kernel vector is
$h_1=(1,0,\ldots,0)$ in flat-character coordinates.  Hence
\[
 \etaInv_1=2\xi_1-h_1
\]
is a $q$-independent invertible translation.  The rank-one twist-eta
reconstruction theorem is known: under one coherently typed multiplicative
character relabeling, equality of the full twist vector forces
$q'\equiv q^{\pm1}\pmod p$
\cite[pp.~23--24, Thm.~2.1; pp.~34--35, Thm.~4.3]{Katase1990}.

\subsection{Adjacent-rank reconstruction is refuted}

Do two adjacent ranks, each retaining the full coherently labelled twist
vector, determine the oriented lens space?  The answer is no.
The all-rank multipliers, chiral reduction, and support algebra above do not
promote to universal oriented reconstruction from two adjacent ranks.  Under
our convention, oriented equivalence of lens spaces requires
$q'\equiv q^{\pm1}\pmod p$, while unoriented homeomorphism allows
$q'\equiv\pm q^{\pm1}\pmod p$; this is the classical topological
classification perimeter \cite{Brody1960}.

The adjacent-rank calculation yields 16 witness rows, all recomputed by both
exact routes in the present convention.  They represent
four pairs of colliding oriented orbits:
\begin{center}
\small
\begin{tabular}{@{}ccllr@{}}
\toprule
$p$ & ranks & left orbit & right orbit & rows\\
\midrule
7  & $(6,7)$ & $\{2,4\}$ & $\{3,5\}$ & 4\\
14 & $(6,7)$ & $\{3,5\}$ & $\{9,11\}$ & 4\\
27 & $(8,9)$ & $\{4,7\}$ & $\{16,22\}$ & 4\\
27 & $(8,9)$ & $\{5,11\}$ & $\{20,23\}$ & 4\\
\bottomrule
\end{tabular}
\end{center}
Every row refutes oriented reconstruction.  The rows at $p=7$ and $p=14$
belong to the blind-rank families proved in Section~\ref{sec:phases}; their
rank vectors carry no information about $q$.  Those eight rows also merge
under orientation reversal, so they are orientation-reversing
homeomorphic.  The eight rows at $p=27$, representing two orbit-pair
collisions, remain nonhomeomorphic even unoriented.  In particular, the
adjacent-rank vectors for $L(27,4)$ and $L(27,22)$ agree exactly with no
character relabeling.

A representative of the second $p=27$ pair is $L(27,5)$ versus $L(27,20)$
with the same relabeling applied to both the reduced vectors and the kernels:
$a\mapsto7a$ (also $a\mapsto20a$).

On the computed domain, the divisibility/size pattern
$(J+1)\mid p$ together with $p\leq3(J+1)$ is an observation, not a theorem.
A refined positive reconstruction criterion remains open.  This negative
result is also distinct from the phase collapse below: $L(16,1)$ and
$L(16,9)$ are exactly distinguished at every computed adjacent pair
$(1,2),\ldots,(8,9)$.  Their all-rank phase loss is therefore not exact-value
loss on this range.

\section{Exponentiated reduced correction and information loss}
\label{sec:phases}

Exponentiation retains the class of the reduced correction modulo integers:
\[
 \tau(B)=\exp(2\pi i\,\xi(B)).
\]
The reduced eta and determinant-line formalism, including the deliberate loss
of the integer lift, is classical
\cite[arXiv v4, pp.~1, 6--8, 12--13]{DaiFreed1994}.  Fractional eta-type information is
related to torsion linking pairings
\cite[arXiv v2, pp.~2, 15, 18]{SavinSternin2002}.  Ordinary Dirac spectra can also lose
geometric information on spherical space forms
\cite[pp.~80--82, Thm.~5]{Bar1996}.  Current second-jet equivariant eta
separations retain different data and concern different lens pairs
\cite[Thm.~1.1, Thm.~4.2, Thm.~5.14, Cor.~5.30]{Sharma2026}.

\subsection{Transform assumptions and Fourier--Dedekind reduction}
Throughout this section, $p\geq2$, $q$ is a unit modulo $p$, the lift
$\widetilde q$ is odd, $u=q^{-1}\pmod p$, $N=J(J+1)$, and $J\geq1$.
We take \eqref{eq:xi-transform} on nonidentity classes and
$\widehat\xi_J(0)=N(2J+1)/12$ as the defining data.  No filling or
unresolved geometric sign choice enters the following proofs.
Set $\alpha=(1+\widetilde q)/2$, $\beta=(1-\widetilde q)/2$ and
\[
 G(x)=x^\alpha\chi_J(x^\alpha)\chi_{J-1}(x^\beta)
     =\sum_e g_e x^e .
\]
The factorization $P_k=x^{-\alpha}(1-x)(1-x^q)$ at $x=\exp(2\pi ik/p)$
gives the exact Fourier inversion
\begin{equation}\label{eq:fd-inversion}
 \xi_{J,a}=\frac{N(2J+1)}{12p}+\sum_e g_e S_{e+a},\qquad
 S_m=\frac1p\sum_{\substack{\lambda^p=1\\\lambda\ne1}}
       \frac{\lambda^m}{(1-\lambda)(1-\lambda^q)}.
\end{equation}
Fourier--Dedekind sums are classical; see \cite{BeckRobins2004}.
Define the ordinary Dedekind sum by
$s(q,p)=\sum_{j=1}^{p-1}(j/p-1/2)((qj\bmod p)/p-1/2)$,
where remainders are least nonnegative integers.  Direct finite summation gives
\begin{align}
 S_m&=\frac1{p^2}\sum_{j=0}^{p-1}j\,((-m-qj)\bmod p)
                    -\frac{(p-1)^2}{4p},\\
 S_0&=-s(q,p)+\frac{p-1}{4p},\qquad
 S_{m+1}-S_m=\frac{(-mu\bmod p)}p-\frac{p-1}{2p}.
\end{align}
For example, the first identity follows from
$(\lambda-1)\sum_{j=0}^{p-1}j\lambda^j=p$ at nonidentity roots and
orthogonality of the $p$th roots.  Summing the difference, also backwards
for negative $m$, yields
\begin{equation}\label{eq:fd-quadratic}
 S_m-S_0\equiv Q(m):=-\frac{u m(m-1)+(p-1)m}{2p}\pmod{\Z}.
\end{equation}

\subsection{The phase law and rank-one generation}
\begin{theorem}[All-rank phase formula]\label{thm:phase-formula}
For the stated transform, every rank and flat label satisfy
\begin{align}
 \xi_{J,a}-\xi_{J,0}&\equiv-\frac{Nu a^2}{2p}\pmod{\Z},
                                      \label{eq:phase-law}\\
 \xi_{J,a}&\equiv-Ns(q,p)-\frac{Nu a^2}{2p}
 -\frac{N(N-2)(\widetilde q+u)}{12p}-\frac N4\pmod{\Z}.
                                      \label{eq:phase-closed}
\end{align}
Both spin lifts give the same phases.
\end{theorem}
\begin{proof}
The first three moments of $G$ are
\begin{align*}
 G_0&:=\sum_e g_e=N,\qquad G_1:=\sum_e eg_e=\alpha N,\\
 G_2&:=\sum_e e^2g_e=\alpha^2N
 +\frac{\alpha^2J^2(J+1)(J+2)+\beta^2J(J+1)^2(J-1)}3.
\end{align*}
The character's first logarithmic derivative at one vanishes; its second
is $\gamma^2r(r+1)(r+2)/3$ for $\chi_r(x^\gamma)$.
Equations~\eqref{eq:fd-inversion}--\eqref{eq:fd-quadratic} give the relative
phase as $-Nu a^2/(2p)$ plus
$-aN(u\widetilde q-1+p)/(2p)$, which is integral.
For the absolute phase, the same equations give
\[
 \xi_{J,0}\equiv\frac{N(2J+1)}{12p}+NS_0
 -\frac{u(G_2-G_1)+(p-1)G_1}{2p}\pmod{\Z}.
\]
Put $t=\widetilde q$ and $ut=1+pk$. Substitution reduces the difference
between this expression and the right side of \eqref{eq:phase-closed}
at $a=0$ to integers: the nontrivial term is
$Nk[t(N+1)+2J+1]/12$.
Indeed $12\mid N[t(N+1)+2J+1]$: modulo three either $3\mid N$ or
$J\equiv1$ and the bracket is zero; modulo four either $4\mid N$ or
$N\equiv2$ and the bracket is even.  Also $N(t-1)/4$ is integral.
Finally $24\mid N(N-2)=(J-1)J(J+1)(J+2)$, so changes of lift or of the
representative of $u$ change the formula by integers.
\end{proof}

The classical congruence
$12p\,s(q,p)\equiv q+u\pmod p$, including the integrality of $12p\,s$,
is an arithmetic property of Dedekind sums \cite[Ch.~3]{RademacherGrosswald1972}.
It simplifies \eqref{eq:phase-closed} to
\begin{equation}\label{eq:phase-simple}
 \xi_{J,a}\equiv-N(N-1)s(q,p)-\frac{Nu a^2}{2p}-\frac N4\pmod{\Z}.
\end{equation}
\begin{corollary}[Rank one generates all phases]\label{cor:rank-one}
With $c_J=J(J+1)/2$ and $\tau_{J,a}=\exp(2\pi i\xi_{J,a})$,
\begin{align}
 \xi_{J,a}&\equiv c_J\xi_{1,a}+2c_J(c_J-1)\xi_{1,0}\pmod{\Z},\\
 \tau_{J,a}&=\tau_{1,a}^{c_J}\tau_{1,0}^{\,2c_J(c_J-1)}.
\end{align}
Thus the complete exponentiated rank ladder carries no more information
than the labelled rank-one phases, under the same coherent relabeling.
\end{corollary}
\begin{proof}
Insert $\xi_{1,a}\equiv-2s-u a^2/p-1/2$ in \eqref{eq:phase-simple}.
The remaining constant difference is $N(N-2)/4\in\Z$.
\end{proof}
This does not identify the full real-valued corrections, nor eliminate the
information in their integer lifts.

\subsection{A complete phase-collapse criterion}
\begin{theorem}[Coherent-label collapse]\label{thm:phase-collapse}
For the same modulus $p$, units $q,q',n$, and any allowed lifts,
$\xi_{J,a}(p,q)\equiv\xi_{J,na}(p,q')\pmod\Z$ for every $J,a$ if and only if
\[
 q'\equiv n^2q\pmod p,\qquad
 2s(q,p)\equiv2s(q',p)\pmod{\Z}.
\]
\end{theorem}
\begin{proof}
At rank one, label zero imposes the Dedekind condition. Subtracting that
case from label one imposes $u\equiv n^2u'\pmod p$, equivalent to the
first condition. Conversely these conditions give equality at every label
of rank one; Corollary~\ref{cor:rank-one} gives every rank.
\end{proof}
The first condition is the usual oriented homotopy relation, weaker than
oriented homeomorphism.  The second is additional arithmetic information.
Equality questions modulo integers for Dedekind sums have an established
literature \cite{Tsukerman2014}; the theorem is an application to the
specified transform, not a new theory of Dedekind sums.

\begin{corollary}[The $p=16$ example]
The pair $L(16,1)$, $L(16,9)$ is not orientation-preservingly homeomorphic,
but has identical phase ladders under $a\mapsto5a$ (also $3a,11a,13a$).
\end{corollary}
\begin{proof}
The oriented homeomorphism criterion fails since $9\ne1^{\pm1}$ modulo 16.
But $5^2\equiv9$ and $12\cdot16\,s(1,16)=210$,
$12\cdot16\,s(9,16)=18$; their difference is $192$, divisible by $6p$.
Apply Theorem~\ref{thm:phase-collapse}.
\end{proof}
The separate quadratic-in-rank calculation
and its 1,536 integer coefficients remain in the supplement as an independent
finite certificate of this example, not the proof of the general theorem.

\subsection{Infinite blind adjacent-rank families}
The following statements concern the exact $\xi$ values, not just phases.
\begin{theorem}[Blind families, including composite moduli]\label{thm:blind}
Let $m\geq5$ be odd and $\gcd(q,m)=\gcd(q^2-1,m)=1$.
For $p=m$ and $r\in\{m-1,m\}$,
\[
 \xi_{r,a}=\frac{r(r+1)(2r+1)}{12p}\qquad\hbox{for all }a.
\]
For $p=2m$, take $q$ odd with the same conditions modulo $m$. Then
\[
 \xi_{r,a}=\frac1p\left\{\frac{r(r+1)(2r+1)}{12}
       +(-1)^{r-1+a}\frac{r(r+1)}4\right\},\quad r=m-1,m.
\]
In either family the two rank vectors are independent of the allowed $q$.
\end{theorem}
\begin{proof}
Writing $z=\zeta^{e_L}$ and $w=\zeta^{e_R}$, centrality means
$p\mid k(1+\widetilde q)$ or $p\mid k(1-\widetilde q)$.
The gcd hypotheses leave only $k=0$ for $p=m$, and $k=0,m$ for $p=2m$.
Since $p\mid2m$, $\chi_{m-1}(z)=0$ whenever $z\ne\pm1$,
and likewise for $w$. Thus the reduced transforms at both ranks vanish
on all noncentral classes. For $p=m$ only the identity term remains.
For $p=2m$, at $k=m$ one has $w=-z$; substitution in the character
polynomials gives $\widehat\xi_r(m)=(-1)^{r-1}r(r+1)/4$.
Fourier inversion multiplies this by $(-1)^a/p$, as asserted.
\end{proof}
For odd $m\geq5$, an admissible $q$ exists exactly when $3\nmid m$.
If $3\mid m$, every unit has $q^2\equiv1\pmod3$, contradicting the
hypothesis. Conversely, $q=2$ works when $3\nmid m$; for the $p=2m$
family, its odd lift $q=m+2$ works. Primality is not necessary: the
hypotheses hold at $m=25,35,49$ with $q=2$, and with $q=3$ at $m=25,35$.
The prime subfamily includes infinitely many nonhomeomorphic collisions:
for primes $p\geq11$, the $p-3$ allowed residues cannot all lie in one
unoriented homeomorphism orbit, whose size is at most four.
The partial-blindness cases $p=3(J+1)$ lie outside these fully blind
families; their classification and a general positive exact-value
reconstruction criterion are not settled by this theorem.

\paragraph{Finite controls, not additional theorem hypotheses.}
The phase formula and rank-one identity each reproduce all 23,688 supplied
table values. The quotient-ring route also checks 3,644 selected prime and
twice-prime entries and 1,690 composite-family entries. These are
reproduced controls, not prospective predictions. Representative-pair
censuses distinguish orientation-preserving from unoriented equivalence;
they must not be read as counts after quotienting every homeomorphism orbit.


Selected exact lifts illustrate the equality modulo integers.
\begin{table}[ht]
\centering
\small
\begin{tabular}{@{}rrrrrr@{}}
\toprule
$J$ & $a$ & $5a\bmod16$ & $\xi_{J,a}(16,1)$ &
$\xi_{J,5a}(16,9)$ & difference\\
\midrule
1 & 1 & 5 & $-7/4$  & $1/4$ & $-2$\\
2 & 1 & 5 & $-7/2$ & $1/2$ & $-4$\\
6 & 1 & 5 & $-1/4$ & $11/4$ & $-3$\\
\bottomrule
\end{tabular}
\caption{Exact phase-lift values for one flat label.  Integer differences
give identical exponentials.}
\label{tab:phase-values}
\end{table}

The kernel term is essential.  The eta-only comparison $[\etaInv/2]$
separates the same pair already at rank two, while the genuine APS class
$[(\etaInv+h)/2]$ does not.  Hence the real eta vector, eta-only reduction,
and exponentiated APS correction retain genuinely different information.
The collapse criterion classifies equality of this exponentiated transform
under coherent relabeling, not homeomorphism of lens spaces.  It does not
classify equality of the full real-valued invariants.

\section{Conclusion and open problems}

The full Dirac completion and its sparse predecessor must be distinguished,
even though their interior symbols are elliptically homotopic. For the
full completion the first-order homogeneous derivation fixes the signed
spectrum and hence the character transform. The auxiliary sector cannot be
discarded: it contributes both
spectral asymmetry and zero modes, and the latter can change what survives
after exponentiation.  At the level of real class transforms the data can be
strongly reconstructive; at the level of exponentiated reduced corrections,
the pair $L(16,1)$ and $L(16,9)$ shows an all-rank loss of oriented
information.  More strongly, the entire phase ladder is generated by rank
one, whereas selected adjacent real-valued ranks can be completely blind to
the lens parameter.  These are distinct information losses.

Several natural questions remain. The fixed-point formula must use the actual
source and target fiber characters \cite[Sec.~4.5]{GilkeyBook1984}; the
character check in Section~\ref{sec:dictionary} is complementary to, not a
replacement for, the first-order operator derivation. The present calculation does not choose a
four-dimensional filling and therefore gives no numerical APS index.  A
positive reconstruction criterion beyond the rank-one theorem is still open,
despite the adjacent-rank counterexamples found here.  It is also unknown how
much of the reducing split and phase law persists for Berger or more general
off-round metrics.  Finally, determinant magnitudes require information not
contained in the reduced correction.  These are analytic and geometric
questions; no physical interpretation is attached to the phase notation used
in this paper.

\section*{Acknowledgment of assistance}
AI-assisted research, computation, checking, and editorial support used
OpenAI Codex and Anthropic Claude. Luqman Mahmood directed the work
and is responsible for the claims and final text. These internal
checks are not external peer review.

\bibliographystyle{alpha}
\bibliography{references}

\end{document}
